% Created 2025-10-13 Mon 08:52
% Intended LaTeX compiler: pdflatex
\documentclass[11pt]{article}
\usepackage[utf8]{inputenc}
\usepackage[T1]{fontenc}
\usepackage{graphicx}
\usepackage{longtable}
\usepackage{wrapfig}
\usepackage{rotating}
\usepackage[normalem]{ulem}
\usepackage{amsmath}
\usepackage{amssymb}
\usepackage{capt-of}
\usepackage{hyperref}
\author{fcalado}
\date{\today}
\title{}
\hypersetup{
 pdfauthor={fcalado},
 pdftitle={},
 pdfkeywords={},
 pdfsubject={},
 pdfcreator={Emacs 29.3 (Org mode 9.6.15)}, 
 pdflang={English}}
\begin{document}

\tableofcontents

\section{draft}
\label{sec:orgffc2e51}
\subsection{introduction}
\label{sec:org196e7fa}
This paper is about normativity, and how the desire for normativity
can connect even the most polarized subjects and experiences.

I take readings of normativity from Trans Studies and apply them to
analysis of cisgendered subjects from a popular reality TV show, \emph{Love
Is Blind}. Participants enter \emph{Love Is Blind} with one goal: to find
their ideal partner and get married. However, the show's main gambit
is that the participants cannot see each other; they must date from
behind a wall, with only access to the other's voice. Once they have
agreed to get married, then they are allowed to meet their intended in
person. And usually, that is when things start to go south. 

My analysis explores what this situation, where visual access to the
beloved is denied to the show's participants, does to the body and to
the self-perception of the body. I explore how this "blind" dating
environment places participants in a state where their own bodily
coherence fractures, and the consequences of this split on the larger
romantic trajectory of the show. I argue that, due to this sensory
dissocation, participants temporarily enter a version of what Jay
Prosser calls the "transsexual trajectory" (4).

Transsexual trajectories, according to Prosser, are characterized by a
desire for "reconcillation between sexed materiality and gendered
identification," and "assimilation, belonging in the body and the
world" (59). As my analysis demonstrates, these dating participants,
while firmly anchored to their cisgendered identities (as far as the
show is concerned, at least), nonetheless experience a kind of split
between what Prosser calls "sexed materiality" and "gendered
identification," in a way that affects their relationship to romance
and marriage, to "belonging in the body and in the world." 

\subsection{toward a new reading practice for normativity}
\label{sec:org3333a29}
To surface these trajectories in the cis subjects of \emph{LiB}, I use a
machine learning method. I use a text generator, which I trained to
"talk like" the show's participants. I chose this method (which I
explain in detail below) because it abstracts from the specific---the
transcripts of the show---to synthesize common and shared investments
in language. It draws from repeated phrasings and frequent words in
the show transcripts in order to generate similar-sounding language.

I also chose this method in attempt to push against dominant uses of
machine learning technology. Amid what Emily Bender and Alex Hanna
call this AI "hype" bubble, it is important to explore versions of AI
that are not environmentally desvasting or extractive as the
commercial-level models which the current tech monopolies are foisting
onto every application they can. The model for this project is a
fraction of the size of the commercial models that run applications
like ChatGPT, meaning the hardware, electricity, and data used for
developing and running the model are miniscule in comparison.\footnote{For example, this project uses GPT-2, which is intially trained
off only 8 million webpages and released under an open license.
Compare that with the most recent version of GPT, GPT-5, which is
trained on something like the entire internet and is over 600 billion
parameters in size, a number that cannot be confirmed due to its
closed and proprietary status. While the training process for my model
took less than an hour, and was all done on my single Mac laptop,
commercial-level models require multiple rounds and massive computer
"clusters," containing hundreds to thousands of GPUs, and months of
training.}

In using this technology as a tool for analysis and critique,
harnessing its affordances for generalizing language in order to study
language use and patterns, I am also pushing back against extractive
logics that prioritize productivity in the promotion of these tools. I
show that this focus on productivity, while serving current capitalist
and imperializing purposes, misses an essential and fascinating aspect
about how these tools work. Within the process of producing langauge,
there are protocols that \emph{normalize} language expression, that find
frequent patterns of word usage, and distill the dominant tendencies
embedded in linguistic aspects.

These tools are, in other words, ideal methods for studying shared
attachments to normativity. For example, when asked about love, the
model will generate examples like:
\begin{quote}
Love is really important to me.

Love is what I need in my life.

Love is everything that I've been looking for.
\end{quote}
None of these examples actually exist in the transcripts of the show.
In fact, in the show transcripts, the only examples of sentences that
begin with "Love is" is, perhaps predictably, followed by the word
"blind." The second half of these sentences, then, are not filled by
verbatim extracts of text of the show: rather, they are filled in by
phrases that appear in similar contexts or are closely associated with
the word "Love" in the transcripts.

Instead of reproducing verbatim expressions, then, the model generates
approximiations and patterns of expressions within the transcripts. In
order to do so, however, it is necessary to "train" the model, so to
speak, to generate text that sounds like the characters on the show.
So that the that the model "learns" how characters talk about certain
topics like love, desire, and the body.

This training process contains several steps, which I will review
briefly. First, a model will randomly guess what words are associated
with each other. In doing so, it takes a sample sentence from the
transcript, like "Love is blind," and it blocks out the second half of
the sentence, so that "Love is" remains. Then, it tries to guess which
word or phrase should go in that blank. Because the model has no idea
what the words mean, the guess will be wrong. But that's doesn't
matter, because the purpose of the hypothesis is to make any guess, so
that it has something from which it build in the future steps.

After making this guess, it moves to the next step, where it checks
its prediction against the actual example. Because the guess will be
inevitably wrong, there will always be a difference between the
prediction and the actual result. And this is when the model starts to
build a sense of the meaning associated with each word, keeping track
of which words tend to appear with each other in context. Each time
that the model makes a subsequent guess, it updates its internal
representation of the word to reflect its association (or lack of
association) to certain words.

With each guess then, the model makes \emph{very slight} adjustments to its
own reprsentation of word meaning. It will make guess after guess,
trying out many words, perhaps every word in the dataset, until it is
sure of those that are most likely to appear together. The model will
do this many times, so progress is very slow, but also very precise.
This constant iteration, and the computer processing required to do
it, is why language models take lots of time, energy, and computer
hardware to train. By the end, the model will have gained a sense of
word meaning by attaining an averaged sense of what that word means in
relation to every other word. 

I read this iterative training process as a kind of \emph{approximation} or
\emph{normalization} of language. The model generates language by
approximating what is most likely, most plausible, based on its
training data. As such, it is ideal for studying shared percpertions
and attachments by the characters on the show.

For example, the phrase "really important," in the actual transcript
of the show, appears in the following excerpt: "I have a great career,
I have a condo, and the one thing I've been missing is finding someone
to share all that with. And that's a really important reason why I'm
here." The model takes this phrase, "really important," and ascertains
what other phrases tend to appear in the same context. It made the
prediction that "love" is associated with concept of lifelong
commitment, and uses that association to generate the text above.

And this is exactly why, while models are good at guessing or
predicting, they are not at all good at being creative, at innovating.
A model can only generate what it has already seen before. Even a
phenomenon like “hallucination,” that a model spews text that has no
bearing in reality, is based on the tendency of models to repeat what
they've already seen. They hallucinate not because they are creative
or random, but because they are designed from statistical processes to
generate what is most plausible rather than most accurate.

\subsection{close reading: love is blind}
\label{sec:org41db745}
In the pods, the body is split: vision is foreclosed, and sensation
is apprehended materially, physically.

In the postpods, vision is restored to the body, there emerges an
aversion to physicality. The aversion to physicality signifies that
characters are not dealing with the reality of their bodies.

The show contains two stages, the pre-engagement stage, where
participants date each other from separate “pods” where they can hear
but not see the other, and engagement stage, where participants
finally meet and proceed to live together in preparation for the
wedding.

In partitioning the romantic experiment into pre-engagement and
engagement segments, the show necessarily poses the presence and role
of the body as the dependent variable that ultimately determines the
viability of long-term commitment. In other words, it sets up
perfectly an examination of how the body may affect normative
trajectories and desires. It allows me to ask two central questions:
first, how does desire emerge when you foreclose the physical body?
sAnd how does this desire change when you introduce the physical body?

To analyze these shifts, I read transcripts from the show. In doing
so, however, I experiment with a new form of close reading that does
not expose, describe, or constellate text—to name some common verbs
used to illustrate close reading analysis. Rather, I use a method that
allows me to study the most normative perspectives expressed by the
characters on the show.

\subsubsection{transgression with a teleology}
\label{sec:org82bfe55}
In the show, participants express their the desire to be "loved for
who I am on the inside," (Season 2, Episode 3). These declarations
recall not only the show's main gambit, that "love is blind," but also
signals a transgressive premise that undergirds the ultimately
conformist plot. Despite the heterosexist teleology, where coupling
ends in lifelong committment, something non-normative is happening
here: a division of the body into parts, the outside and the inside,
the parts that we can see, and the parts that we can feel.

This places the characters into a trajectory that is not unlike what
Jay Prosser calls "transsexual narrative," at least up to a point.
According to Jay Prosser, "transsexual narratives" are marked by
teleology, in which "transition may be the very route to identity and
bodily integrity" (Prosser 6). I argue below that the characters,
although they are cis subjects, undergo a process that is something
like transition in the dating experience of the pods. In being severed
from visually seeing each other, they exprience a kind of sensory
dissociation, which inevitably places them, along with the marriage
plot, in a journey toward "bodily integrity," in Prosser's words.

\subsubsection{touch into the body and out of the body}
\label{sec:orgb53bcc1}
Initial signs of this journey emerge in the contrast in how characters
approach physical touch during their time in the pods, when they are
physically separated from each other, and in the time following their
meeting, when they live together for several weeks prior to the
wedding.

To explore the differences between these periods of the show, I
trained two separate models on each of these periods: the first model
is trained on the transcripts from the pods episodes, and I call it
the "pods model"; the second model is trained on the episodes
following the pods, and I call it the "postpods model."

Then, I prompted both models with input phrases like "When I touch
you," and "Physical touch is," and "Physically," (See Appendices 1 \& 2
for a full list of prompts), to see how each of them would respond
individually to the prompts. 

In the pods model, touch is mostly associated with abstract or
non-tangible aspects of the other:
\begin{quote}
When I touch you, I feel it.

When I touch you, I feel your energy, and it is, like, I'm just so
happy.

When I touch you, I can feel your soul, your heart, and your soul
aligning so well.

When I touch you, I feel you, and I feel you, I feel you, and I feel
you.

Physical touch is the most important thing.

Physical touch is important to me.

Physical touch is so sexy.

Physical touch is like a glove.

Physically, we are so happy.
\end{quote}
Considering that all these descriptions occur when there is no actual
touching occuring between the couples, it makes sense that the model
would associate touch to non-tangible aspects like "soul" and
"energy," to things that have no physical existence (at least in a
scientifically verifiable sense). It also makes sense that touch is
elevated to an "important" and even "sexy" quality, because that is
what is being foreclosed from during this period of the dating
experiment, and therefore represents something desired. The last two
examples, "like a glove" and "we are so happy," signal compatability,
in the sense that they fit and are happy together.  

The outputs with those from the postpods model, however, put touch in
a very context:
\begin{quote}
When I touch you, I feel like I'm in my head.

When I touch you, you just feel like it's so weird.

When I touch you, it feels like a jab.

When I touch you, like, I feel like I'm literally in my head.

When I touch you, it feels like something I'm about to get up and
walk away.

When I touch you, I feel like it's like I've just, like, left the
room.

When I touch you, the thing that's scary is, like, it's a physical
thing.

When I touch you, you're like "I'm blinking."
\end{quote}
Unlike the pods model, these descriptions present touch as a not
pleasant experience. Touch is "so weird," "like a jab," and associated
with workds like "scary" and "physical," suggesting a strange
surprising, and even disruptive aspect to touch. Touch also leads to
dissociating from physical location, describing movements from "in my
head" to "out of the room." While in the pods, touch drew the
characters together, making evoking and sensualizing non-tangible
phenomena like the soul and energy, here it seems that touch pulls
characters away from their own bodies and each other.

\subsubsection{"blinking" out}
\label{sec:org41a0230}
The peculiar reference to "blinking" in the last output recalls one
particular narrative thread, that between the characters Zach and
Irina. The phrase, "When I touch you, you're like 'I'm blinking,'"
actually results from a blip in the algorithmic prediction process. It
is contains a quote taken directly from the show. In machine learning,
it referred to as "overfitting," when a verbatim section of text from
the training data, in this case, the show transcripts, is generated in
the output. "Overfitting" means that the model is \emph{too} accurate: that
it has slipped from making predictions to parroting the data it has
been trained on. A model overfitting in its outputs is generally a
sign that there isn't enough training data. With less training data,
there is less variation, meaning that the model has less examples from
which to approximate and generalize. So it results it simply
reproducing direct examples from its training. 

For my purposes, however, overfitting is not only a blip, it also
opens a new dimension for understanding "touch." It leads to specific
scene in the text, a first meeting between a newly engaged couple,
Zach and Irina: 
\begin{quote}
Zach: Do I look like what you thought I'd look like?

Irina: I had no guesses of what you looked like.

Zach: Oh!

Irina: You have, like, the blankest stare in your eyes.

Zach: Really?

Irina: I'm just kind of taking it all in.

Zach: Me too.

Irina: You look like a fictional character. You look like something
out of a cartoon.

Zach: I know.

Irina: You have to blink!

Zach: I am blinking.

Irina: You don't blink. You look like this.

Zach: I am blinking. I will try not to be too intense. \emph{Season 4,
Episode 4, "Playing with Fire"}
\end{quote}
When Zach and Irina meet, it is clear that both are adjusting to the
reality of each other's physical form in their own ways. Irina, in
particular, appears overwhelmed to the point of asking him to change
his behavior, "you have to blink!" Blinking is, of course, a way of
stopping the entry of visual data, of occluding it from perception.
For Irina, the request for Zach to blink indicates her own sense of
overwhelm at his physical form, his sudden incorporation. The visual
reality of him takes her by surprise, even overstimulates her, perhaps
leading her to project that sense onto him, which is why she asks him
to blink.

Following the story of Zach and Irina, it becomes clear that one of
the reasons for their breakup is a lack of physical attraction on the
part of Irina. Later in the same episode, Irina explains her feelings
to another character, Micah, who is coupled with Paul.
\begin{quote}
Irina: And so, Zack. I feel like is my type on paper. Has, like, brown
hair, brown eyes, like, chiseled face. Like, I really like dark
features. And the moment I saw Zack, it was like, "I don't know who
this man is." And I was like, "Maybe it's just scary, and it was a
lot." Like, hopefully it's gonna grow, but I've noticed every time he
does, like, touch me, I get, like, major ick. When he puts his arm
around me at night, I literally was like-- like, my heart stopped. And
I literally go\ldots{} But not, like, in an excited way.

Micah: I wanna, like, relate to you in a way, but it's always, like,
so different.

Irina: How was it with you and Paul?

Micah: The thing with me and Paul is, like, we both, like, had such an
immediate understanding as best friends.

Irina: Yeah, Paul's gorgeous.
\end{quote}
Zach supposedly has physical aspects which Irina finds attractive,
"brown eyes, chiseled face," but something about him nonetheless
repulses her. At his touch, she recoils, "get[ting]\ldots{} major ick."
While she claims her attraction is not determined by a visual sense,
Irina, with no particular prompting, elevates visual appearance to the
criterion for attraction with the phrase, "Yeah, Paul's gorgeous."
Which begs the question: if Irina finds both Zach and Paul visually
attractive, then, what repulses her about Zach?

The answer could be a number of things, which Irina may or may not
have revealed in the course of the show. But the point is not the
specific reason. Rather, it's that the model made this connection
between touch and "blinking." This connection, then, posits that the
visual sense can overwhelm and occlude the other senses, such as
touch, which is a result of the experience of the characters in the
pods. There is something about being the pods, about the foreclosure
of the visual which activates a sensual experience of embodiment that
is subsequently ruptured once the characters meet in person.

\subsubsection{my body is coming off}
\label{sec:org115d49c}
Within their experience in the pods, these cisgendered characters
experience their body in new ways. They experience not only the
physical body, the material reality of their physical body, but what
Jay Prosser refers to as the "body image," a separate but nonetheless
material perception of the body that exists within the material body.
According to Prosser, "body image---which we might be tempted to align
with the imaginary---clearly has a material force for transsexuals"
(69). This "material force" can manifest in many ways, and often does
so in the trope of being "trapped in the wrong body" and feelings of
dysphoria. For trans subjects, the body image creates "conflict
between the true body within and the false body without, between
sentient body image and insentient visible body" (70).

The pods environment, I argue, faciliates the emergence of embodied
sensation that has a new and suprising force for these cisgendered
characters. Below is the results for the pods model when prompted with
the phrase "My body":
\begin{quote}
My body makes me feel like it's real.

My body feels weird.

My body feels heavier.

My body feels so different now.

My body feels like it's coming off.

My body feels torn between two different people.
\end{quote}
These outputs describe an awareness of the body coming into
apprehension in a novel and visceral way, an awareness which makes the
body seem all the more strange. The body feels "weird" and "so
different now" because it is not the same body as that which also has
access to the visual sense. Coming into perception, this new body is
"heavier," more weighty, more material. At the moment where visual
access to the other's body is denied, there is an increased sensing of
one's own physicality. And this heightened sense of body image is not
fill a sensual lack for that which has been foreclosed, the sight of
the other's body. It is not that one sense compensates for another.
Rather, it is that the foreclosure of sight bifurcates the material
body, which was previously whole, into the body image. Which is why,
perhaps, "The body feels torn between two different people."

Perhaps, for the straight cisgendered folks within the pods, this is
the first time they've felt the divergence between the material and
imagined bodies. Despite being cisgendered, then, body image can have
a material force in environment where the integrity bodily perception
is disrupted. The sensation of the body feeling "real" and like it's
"coming off" speak to this process of fragmentation, a process that is
so visceral so that, as Prosser explains, "feeling like\ldots{} can be
experienced as a core self" (79).

The pods facilitate an experience of embodied sensuality that is, for
most of the characters, subsequently foreclosed when they leave the
pods. After the couples meet in person, descriptions of the body shift
more firmly into visual.

\begin{quote}
My body is gorgeous.

My body is so cute.

My body is so pretty.

My body makes me feel lighter, more confident.

My body makes me feel warm.

My body makes me feel like I've missed my train.
\end{quote}
The outputs address the body's physical appearance in positive terms:
the body is "gorgeous," "so cute," "pretty" (Appendix 2, "My body").
Descriptive terms refelct that now the progression of the relationship
to in-person, where the couples can now see the other. This visual
experience leads to a kind of integration: it is now lighter" and
"warm," suggesting a bodily coherence where before there was weight
and weirdness. Despite this coherence, however, there is a feeling of
loss: "My body makes me feel like I've missed my train." This
statement, which diverges a bit from the others, leaves a trace or a
doubt that even when coherence is gained, there is something lost.

\subsubsection{vision restored, something is lost}
\label{sec:org52939ea}
When prompted on physical touch, the sense of what is lost comes into
view. It is a sense of potential, perhaps resulting from the
incoherence of the body.
\begin{quote}
Physical touch is everything that I've ever wanted in a partner.
Physical touch is everything that I've wanted in a wife.
Physical touch is a big part of what I want.
Physically, there's so much potential here.
Physically, it was the perfect opportunity.
\end{quote}
Here, physical touch is described in aspirational terms: it is
"everything I've wanted," not "everything I got." The traces of
present perfect and past tense to discuss desire is indicative here:
even after meeting in person, their desire remains unfulfilled. The
restoration of the visual sense, the integration the previously
fractured body image does not offer completion or culmination.

Being restored their visual sense heals them from the bodily split,
but it does not save them from the aftermath of their experience in
the pods. Although they exited from one trajectory, what Prosser calls
the "route to bodily integrity," they remain stuck at the end with
desires for a kind of touch that is "everything that I've ever wanted
in a partner." Considering that the characters are now reunited with
their physical bodies, there is something cruel in this denouement, a
"cruel optimism," in Lauren Berlant's formulation, which describes the
attachment that drives desire even while it wears out the desirer.

They experience this desire because, even before the bodily split,
subjects had formed attachments to a romantic love that trascends
physical barriers. Then, when they enter the pods, the foreclosure of
one sense perception instigates and amplifies embodied sense. This new
body image emerges seemingly from nowhere, and feeds them sensory
delights from behind the wall. Subsequently, when the couples finally
meet in physical forms, they remain plagued by the possibilities for
physical connection that they felt in the pods. And these expectations
are what, for some of them, prevents their ability to accept their
partners as they are. 

More than cruel optimism, then, the characters perhaps experience
something closer to what Hil Malatino describes as "future fatigue,"
in which anxiety builds for a dream that is ultimately deferred. Like
cruel optimism, future fatigue generates "intense anticipatory
anxiety" that "impede[s] flourishing" (Malatino 20). Unlike cruel
optimism, however, future fatigue specifically concerns trans subjects
who are invested in "achieving the promised moment of harmony between
the felt and the perceived body" (Malatino 27). Despite being firmly
cisgendered (whatever that means), these subjects experience a
disruption between the felt and the perceived bodies.

\subsubsection{toward a kind of erotics}
\label{sec:orgeb391af}
Malatino theorizes "future fatigue" as a foundational affect from
which to read moments of potential and connection among trans
subjects. Malatino's project explores how critical attention might
shift its reading of frustration and failure toward recognizing
"space[s] of looseness and possbility, not yet overcoded and fixed in
meaning, signification, or representative economy" (32). I am not
interested in reading, as Malatino does, moments of potential and
connection among these characters. Cis-hetero significations and
dynamics are well respresented and coded. What I am interested in,
rather, is in surfacing some vectors of connection between cis and
trans experience. Because I believe that this approach could help push
back on the polarization that characterizes the current discourse
about trans identity and rights in the USA.

And I know this because I began this study from that polarized point
of view. When I first started this project, my goal was to study
transphobia directly. To use generated text to study fear. I gathered
datasets from conservative magazines and anti-trans legislation, and
used those to train my machine learning models. I thought machine
learning, with its tendency to amplify what is most frequent, would be
an ideal tool for surfacing perspectives around transphobia.

As I got deeper and deeper into transphobia, however, I realized that
perhaps I'm working with the wrong data. That I should have been
studying attachment to norms, particularly the attachments to sexual
and romantic norms. That route might be a way to bring cis and trans
more firmly into relation. Although this might be a controversial
question, I think it's a necessary one in our current political
moment.

I’ll finish with a final question by Cassius Adair. In his study for
t4t erotics, he asks, "Why can't the erotic be a site of producing
trans identity or practices?" He points out that, after all, cis
people do it all the time: they use sexuality and sexual encounters as
sites of identity formation."
\end{document}
